\chapter{Canonical Transformations, Generators, and Symmetry}

Phase space is not merely a list of positions and momenta. It has a structure,
and the structure is what remains meaningful when the coordinates used to
describe it are changed. The Poisson bracket is the local expression of that
structure. It tells us which changes of phase-space coordinates are allowed,
how functions flow under those changes, and why a symmetry carries a conserved
quantity with it.

\section{Structure and Flow on Phase Space}

There are many kinds of structure one can place on a space. A Riemannian
metric, for example, supplies a distance between neighboring points. A
Poisson structure is different: it supplies an antisymmetric bracket between
functions. For mechanics the coordinates of the space are not only
\(q_1,\ldots,q_n\), but the full set
\[
 (q_1,\ldots,q_n,p_1,\ldots,p_n).
\]
The question is therefore richer than a change of coordinates in ordinary
space. We want to know which transformations may mix the \(q\)'s and \(p\)'s
without destroying the form of mechanics.

It is useful to think first about familiar transformations. A finite rotation
is built from many small rotations, and a finite translation from many small
translations. The intermediate configurations trace a flow. The flow need
not be the actual motion of a physical system; it can simply carry one
description into another.

The same idea applies to phase space. Hamiltonian evolution is one flow: a
phase point moves with time under the influence of \(H\). Rotations,
translations, and less familiar transformations that mix positions and
momenta define other flows. Poisson brackets give one language for all of
them. Newton's equations are the most concrete end of this story; the Poisson
formulation is the abstract end needed here.

Ordinary mechanical systems most visibly exhibit translations and rotations,
although special systems can hide additional symmetries. More advanced
physics adds internal symmetries that are not motions of ordinary space at
all. The common question is structural: what can be transformed while the
laws retain their form? Hamilton and Poisson supplied the classical language
in which that question can be asked directly on phase space.

\section{The Poisson Algebra}

Let \(A\), \(B\), and \(C\) be functions on phase space. The first defining
property is antisymmetry,
\begin{equation}
 \{A,B\}=-\{B,A\}.
 \label{eq:l09-antisymmetry}
\end{equation}
The second is linearity in each entry. For ordinary numbers
\(\alpha\) and \(\beta\),
\begin{equation}
 \{\alpha A+\beta B,C\}
 =\alpha\{A,C\}+\beta\{B,C\}.
 \label{eq:l09-linearity}
\end{equation}

\begin{figure}[tbp]
\centering
\includegraphics[width=0.76\textwidth]{lecture_09_figure_01.png}
\caption{Antisymmetry and linearity on the board at 00:09:12. The minus sign
in the first rule is the feature that fixes the order of every later bracket.}
\label{fig:l09-axioms-board}
\end{figure}

\begin{classroomqa}
\audiencequestion{Should the first relation have a minus sign when \(A\) and
\(B\) are interchanged?\lecturetimestamp{00:08:09--00:08:18}}
\lecturerresponse{Yes. That minus sign is the essential point:
\(\{A,B\}=-\{B,A\}\). Without it there would be no distinction between the
two orders.}
\end{classroomqa}

The bracket also obeys the product rule,
\begin{equation}
 \{AB,C\}=A\{B,C\}+B\{A,C\}.
 \label{eq:l09-product-rule}
\end{equation}
By antisymmetry the corresponding rule in the second entry follows as well.

\begin{figure}[tbp]
\centering
\includegraphics[width=0.76\textwidth]{lecture_09_figure_02.png}
\caption{The product rule at 00:11:18. In classical mechanics the functions
outside the bracket multiply as ordinary commuting numbers.}
\label{fig:l09-product-board}
\end{figure}

The order of multiplication outside a bracket is immaterial:
\(A\{B,C\}=\{B,C\}A\). The order \emph{inside} the bracket is not
immaterial, because reversing it changes the sign. This distinction is easy
to overlook in classical mechanics, where functions commute, but it
anticipates the ordering questions that become unavoidable in quantum
mechanics.

\begin{classroomqa}
\audiencequestion{What sort of objects are \(A\) and \(B\)? Must they be
real-valued scalars?\lecturetimestamp{00:09:26--00:14:13}}
\lecturerresponse{They are numerical-valued functions of all the phase-space
coordinates. They may be added and multiplied pointwise, and complex-valued
functions are allowed when useful. Here \emph{numerical} only means that their
ordinary multiplication is commutative; it is separate from the physicist's
use of \emph{scalar} to mean invariant under spatial rotations.}
\end{classroomqa}

For canonical coordinates the concrete derivative formula is
\begin{equation}
 \{A,B\}
 =\sum_{i=1}^{n}\left(
 \frac{\partial A}{\partial q_i}\frac{\partial B}{\partial p_i}
 -\frac{\partial A}{\partial p_i}\frac{\partial B}{\partial q_i}
 \right).
 \label{eq:l09-poisson-definition}
\end{equation}
We assume throughout that the functions are smooth enough for the indicated
derivatives and limiting arguments.

\section{Canonical Brackets as Derivatives}

Antisymmetry immediately gives every self-bracket:
\[
 \{p_i,p_i\}=0,\qquad \{q_i,q_i\}=0.
\]
The derivative definition gives the complete canonical relations,
\begin{equation}
 \{q_i,q_j\}=0,\qquad
 \{p_i,p_j\}=0,\qquad
 \{q_i,p_j\}=\delta_{ij}.
 \label{eq:l09-fundamental-brackets}
\end{equation}
Consequently \(\{p_i,q_j\}=-\delta_{ij}\). The Kronecker delta is one
when \(i=j\) and zero otherwise.

\begin{figure}[tbp]
\centering
\includegraphics[width=0.74\textwidth]{lecture_09_figure_03.png}
\caption{The canonical brackets at 00:17:25. The mixed bracket is the
Kronecker delta; brackets among two coordinates or two momenta vanish.}
\label{fig:l09-canonical-board}
\end{figure}

Antisymmetry determines a self-bracket because it equals its own negative.
For two different variables it says only
\(\{q,p\}=-\{p,q\}\); the normalization \(\{q,p\}=1\) comes from the
canonical derivative definition.

Two useful zeroes follow at once:
\begin{equation}
 \{F(q),K(q)\}=0,\qquad \{G(p),L(p)\}=0.
 \label{eq:l09-same-type-zero}
\end{equation}
Each term in a Poisson bracket contains one derivative with respect to a
coordinate and one with respect to a momentum. A nonzero contribution
therefore needs a coordinate dependence on one side to meet the conjugate
momentum dependence on the other.

The fundamental brackets and the product rule turn bracketing into
differentiation. For one canonical pair,
\begin{equation}
 \{F(q),p\}=\frac{dF}{dq}.
 \label{eq:l09-fq-p}
\end{equation}
Rather than simply reading this from
Eq.~\eqref{eq:l09-poisson-definition}, let us recover it from the algebra.

For the constant function,
\begin{align}
 \{1^2,p\}
 &=1\{1,p\}+\{1,p\}1
 =2\{1,p\}.
\end{align}
Since \(1^2=1\), this says \(\{1,p\}=2\{1,p\}\), hence
\(\{1,p\}=0\). Next,
\[
 \{q,p\}=1.
\]
For the next monomial,
\begin{equation}
 \{q^2,p\}
 =q\{q,p\}+\{q,p\}q
 =2q.
 \label{eq:l09-q2-proof}
\end{equation}

\begin{figure}[tbp]
\centering
\includegraphics[width=0.70\textwidth]{lecture_09_figure_04.png}
\caption{The monomial proof at 00:22:43: the cases \(1\), \(q\), and
\(q^2\) reproduce their ordinary derivatives.}
\label{fig:l09-monomial-board}
\end{figure}

Induction gives \(\{q^r,p\}=r q^{r-1}\). Linearity extends the result to
polynomials and, under the smoothness assumption, to the functions used in
mechanics. In many variables the paired identities hold for a general
phase-space function:
\begin{equation}
 \{F(\bm q,\bm p),p_i\}=\frac{\partial F}{\partial q_i},
 \qquad
 \{q_i,F(\bm q,\bm p)\}=\frac{\partial F}{\partial p_i}.
 \label{eq:l09-bracket-derivatives}
\end{equation}
The order is deliberate: \(p_i\) goes on the right to differentiate with
respect to \(q_i\), whereas \(q_i\) goes on the left to differentiate with
respect to \(p_i\).

\begin{classroomqa}
\audiencequestion{Does the extension from polynomials require continuity of
the bracket and of the functions?\lecturetimestamp{00:24:11--00:24:38}}
\lecturerresponse{Yes. We are assuming smooth, multiply differentiable
functions and the usual continuity needed to pass from polynomial
approximations to their limits.}
\end{classroomqa}

\section{Hamiltonian Evolution Is One Flow}

Dynamics enters through one further relation:
\begin{equation}
 \dot A=\{A,H\}.
 \label{eq:l09-hamiltonian-evolution}
\end{equation}
A phase point follows an orbit determined by \(H\). As the point moves, a
function \(A(q,p)\) changes along that orbit, and
Eq.~\eqref{eq:l09-hamiltonian-evolution} is its directional derivative along
the Hamiltonian flow.

\begin{classroomqa}
\audiencequestion{From the abstract viewpoint, are the \(p\)'s and \(q\)'s
being assumed to depend on time?\lecturetimestamp{00:27:20--00:27:52}}
\lecturerresponse{Apply the same rule to the coordinates themselves:
\(\dot q_i=\{q_i,H\}\) and \(\dot p_i=\{p_i,H\}\). Their time dependence is
then defined by the Hamiltonian flow.}
\end{classroomqa}

For a free particle in one dimension,
\begin{equation}
 H=\frac{p^2}{2m}.
\end{equation}
The momentum equation is
\begin{equation}
 \dot p=\left\{p,\frac{p^2}{2m}\right\}=0,
\end{equation}
because both entries depend only on \(p\). The coordinate equation is
\begin{align}
 \dot q
 &=\left\{q,\frac{p^2}{2m}\right\} \notag\\
 &=\frac{1}{2m}\left(p\{q,p\}+\{q,p\}p\right)
 =\frac{p}{m}.
 \label{eq:l09-free-flow}
\end{align}

\begin{figure}[tbp]
\centering
\includegraphics[width=0.75\textwidth]{lecture_09_figure_05.png}
\caption{The free-particle calculation at 00:29:22. The Poisson algebra gives
\(\dot p=0\) and \(\dot q=p/m\) directly.}
\label{fig:l09-free-flow-board}
\end{figure}

Thus the abstract bracket algebra reproduces Hamilton's equations and the
ordinary velocity--momentum relation.

\begin{classroomqa}
\audiencequestion{Is this Poisson formulation restricted to nonrelativistic
mechanics?\lecturetimestamp{00:29:57--00:30:45}}
\lecturerresponse{No. The example \(H=p^2/(2m)\) is nonrelativistic, but the
canonical structure applies much more broadly: relativistic particles,
classical fields, electromagnetism, and general relativity all admit the
appropriate Hamiltonian form. Quantum mechanics replaces the classical
bracket by its operator counterpart.}
\end{classroomqa}

\section{Which Coordinate Changes Preserve Mechanics?}

We now separate two questions. First, which changes of phase-space
coordinates preserve the Poisson structure? Second, among those changes,
which leave a particular Hamiltonian invariant? The first question defines a
canonical transformation; the second will define a symmetry.

Take one canonical pair \((q,p)\). A simultaneous stretch,
\begin{equation}
 Q=2q,\qquad P=2p,
\end{equation}
fails because
\[
 \{Q,P\}=4\{q,p\}=4.
\]
A compensated stretch and squeeze,
\begin{equation}
 Q=2q,\qquad P=\frac12 p,
 \label{eq:l09-squeeze}
\end{equation}
succeeds:
\[
 \{Q,P\}=\left\{2q,\frac12p\right\}=1.
\]

\begin{figure}[tbp]
\centering
\includegraphics[width=0.72\textwidth]{lecture_09_figure_06.png}
\caption{The compensated rescaling at 00:33:55. Stretching \(q\) while
squeezing \(p\) preserves the canonical bracket.}
\label{fig:l09-squeeze-board}
\end{figure}

For one pair, the successful transformation preserves area in the
\((q,p)\)-plane. The double stretch does not. In many dimensions, preserving
total volume is not enough: the transformation must preserve the full
symplectic form
\[
 \omega=\sum_i dq_i\wedge dp_i,
\]
or equivalently the complete canonical bracket relations.

\begin{classroomqa}
\audiencequestion{Is the squeeze-and-stretch condition basically conservation
of area?\lecturetimestamp{00:34:51--00:35:20}}
\lecturerresponse{For one \(q,p\) pair, yes: canonical transformations are
area-preserving maps of that plane. With several pairs, ordinary volume
preservation is only a consequence; preservation of the full symplectic
structure is the stronger condition.}
\end{classroomqa}

A rotation of the phase-space plane also works:
\begin{equation}
 P=\cos\theta\,p+\sin\theta\,q,\qquad
 Q=-\sin\theta\,p+\cos\theta\,q.
 \label{eq:l09-phase-rotation}
\end{equation}

\begin{figure}[tbp]
\centering
\includegraphics[width=0.74\textwidth]{lecture_09_figure_07.png}
\caption{A phase-space rotation at 00:36:10. Positions and momenta are mixed,
but the canonical structure survives.}
\label{fig:l09-rotation-board}
\end{figure}

Self-brackets vanish by antisymmetry, so only the mixed bracket needs work:
\begin{align}
 \{Q,P\}
 &=\{-\sin\theta\,p+\cos\theta\,q,\,
       \cos\theta\,p+\sin\theta\,q\}\notag\\
 &=\cos^2\theta\,\{q,p\}
   -\sin^2\theta\,\{p,q\}\notag\\
 &=\cos^2\theta+\sin^2\theta=1.
 \label{eq:l09-rotation-check}
\end{align}

\begin{definition}
A transformation
\[
 (q_i,p_i)\longmapsto(Q_i,P_i)
\]
is canonical when
\[
 \{Q_i,Q_j\}=0,\qquad
 \{P_i,P_j\}=0,\qquad
 \{Q_i,P_j\}=\delta_{ij}.
\]
\end{definition}

\begin{classroomqa}
\audiencequestion{Are the equations of motion unchanged after a canonical
transformation?\lecturetimestamp{00:40:04--00:40:20}}
\lecturerresponse{Their canonical form is unchanged, but the Hamiltonian must
be re-expressed as a function of the new variables \(Q\) and \(P\). The
numerical formula for \(H\) need not look the same.}
\end{classroomqa}

\begin{classroomqa}
\audiencequestion{Are canonical transformations examples of gauge
transformations?\lecturetimestamp{00:40:21--00:41:05}}
\lecturerresponse{No. A gauge transformation relates redundant
descriptions of the same physical state. A canonical transformation is an
ordinary change of coordinates on phase space that preserves its Poisson
structure.}
\end{classroomqa}

\begin{classroomqa}
\audiencequestion{Can a canonical transformation carry one arbitrary
physical system into another?\lecturetimestamp{00:41:42--00:44:52}}
\lecturerresponse{Here it means a change of phase-space coordinates for the
same system. Phase spaces with the same number of degrees of freedom may be
locally isomorphic, but a canonical re-description does not by itself turn a
harmonic oscillator into an unrelated many-body system.}
\end{classroomqa}

\section{Infinitesimal Transformations and Their Generators}

\begin{classroomqa}
\audiencequestion{Is there a more constructive test than checking all the
new Poisson brackets one by one?\lecturetimestamp{00:45:07--00:46:31}}
\lecturerresponse{For an infinitesimal transformation there is. Build the
change from a generator \(G(q,p)\). Finite transformations connected to the
identity can then be assembled from these small steps.}
\end{classroomqa}

Write one small step as
\begin{equation}
 Q=q+\delta q,\qquad P=p+\delta p,
 \label{eq:l09-infinitesimal-change}
\end{equation}
where \(\delta q\) and \(\delta p\) are themselves functions of phase space.
Infinitesimal means that terms quadratic and higher in these changes may be
discarded. Expanding the transformed bracket gives
\begin{align}
 \{Q,P\}
 &=\{q+\delta q,p+\delta p\}\notag\\
 &=\{q,p\}+\{\delta q,p\}+\{q,\delta p\}
   +\{\delta q,\delta p\}.
 \label{eq:l09-first-order-expansion}
\end{align}
The last term is second order. Therefore the first-order canonical condition
is
\begin{equation}
 \{\delta q,p\}=-\{q,\delta p\}.
 \label{eq:l09-first-order-condition}
\end{equation}

\begin{figure}[tbp]
\centering
\includegraphics[width=0.78\textwidth]{lecture_09_figure_08.png}
\caption{The first-order canonical condition at 00:50:44:
\(\{\delta q,p\}=-\{q,\delta p\}\).}
\label{fig:l09-first-order-board}
\end{figure}

Introduce a small parameter \(\epsilon\) and a function \(G(q,p)\). Define
\begin{equation}
 \delta q=\epsilon\{q,G\},\qquad
 \delta p=\epsilon\{p,G\}.
 \label{eq:l09-generator-definition}
\end{equation}
Using Eq.~\eqref{eq:l09-bracket-derivatives},
\begin{equation}
 \delta q=\epsilon\frac{\partial G}{\partial p},
 \qquad
 \delta p=-\epsilon\frac{\partial G}{\partial q}.
 \label{eq:l09-generator-derivatives}
\end{equation}

\begin{figure}[tbp]
\centering
\includegraphics[width=0.78\textwidth]{lecture_09_figure_09.png}
\caption{The generator equations at 00:53:55. A single function \(G\)
determines both components of the infinitesimal phase-space flow.}
\label{fig:l09-generator-board}
\end{figure}

These formulas solve the canonical condition automatically:
\begin{align}
 \{\delta q,p\}
 &=\epsilon\left\{\frac{\partial G}{\partial p},p\right\}
 =\epsilon\frac{\partial^2G}{\partial q\,\partial p},\notag\\
 -\{q,\delta p\}
 &=\epsilon\left\{q,\frac{\partial G}{\partial q}\right\}
 =\epsilon\frac{\partial^2G}{\partial p\,\partial q}.
 \label{eq:l09-mixed-derivative-proof}
\end{align}
The equality is simply equality of mixed partial derivatives.

\begin{figure}[tbp]
\centering
\includegraphics[width=0.72\textwidth]{lecture_09_figure_10.png}
\caption{The mixed-derivative check at 00:55:50. It proves that a
generator-defined infinitesimal flow is canonical.}
\label{fig:l09-mixed-proof-board}
\end{figure}

\begin{classroomqa}
\audiencequestion{Must the two variations use the bracket with \(G\) in the
same order?\lecturetimestamp{00:52:08--00:52:15}}
\lecturerresponse{Yes:
\(\delta q=\epsilon\{q,G\}\) and
\(\delta p=\epsilon\{p,G\}\). The opposite signs in their derivative forms
come from the canonical bracket itself, not from reversing the order.}
\end{classroomqa}

The same generator controls every function. By the chain rule,
\begin{align}
 \delta A
 &=\frac{\partial A}{\partial q}\delta q
   +\frac{\partial A}{\partial p}\delta p\notag\\
 &=\epsilon\left(
 \frac{\partial A}{\partial q}\frac{\partial G}{\partial p}
 -\frac{\partial A}{\partial p}\frac{\partial G}{\partial q}
 \right)\notag\\
 &=\epsilon\{A,G\}.
 \label{eq:l09-general-generated-flow}
\end{align}
For many degrees of freedom the sum over \(i\) is understood, and each
\(\delta q_i\) may depend on every \(q_j\) and \(p_j\), not only on the pair
with the same index.

Hamiltonian evolution is the special choice
\[
 G=H,\qquad \epsilon=\delta t.
\]
Thus time evolution is itself a canonical transformation. More generally,
finite canonical transformations in the component connected to the identity
are built by composing infinitesimal generator flows. This is the precise
local statement needed here; global topology can introduce qualifications.

\begin{classroomqa}
\audiencequestion{If \(G\) generates a flow in the same way as \(H\), is it
just the Hamiltonian of another system?\lecturetimestamp{00:58:35--00:59:14}}
\lecturerresponse{Mathematically it may be treated that way. Any suitable
\(G\) defines a Hamiltonian-like flow. Which function is the actual energy
of the physical system is additional physical information, fixed by the
system and ultimately by experiment.}
\end{classroomqa}

\section{Symmetry and Conservation}

A canonical transformation preserves the form of mechanics. A symmetry of a
particular system must do more: it must preserve that system's Hamiltonian.
Suppose \(H\) generates the physical time flow and \(G\) generates another
canonical flow. In general the \(G\)-flow crosses surfaces of constant
energy. It is a symmetry precisely when it runs along them.

\begin{figure}[tbp]
\centering
\includegraphics[width=0.72\textwidth]{lecture_09_figure_11.png}
\caption{The phase-space flow picture at 01:02:10. The curved trajectories
represent Hamiltonian motion; a second flow is a symmetry when it remains on
the same energy surface.}
\label{fig:l09-flow-board}
\end{figure}

\begin{figure}[tbp]
\centering
\resizebox{\linewidth}{!}{%
\begin{tikzpicture}[x=1.15cm,y=1.0cm,>=Latex]
  \draw[->] (0,0) -- (6.4,0) node[right]{phase-space direction};
  \draw[->] (0,0) -- (0,4.2);
  \foreach \y/\lab in {0.8/{H=E_1},2.0/{H=E_2},3.2/{H=E_3}} {
    \draw[thick]
      plot[smooth] coordinates {(0.5,\y) (2.0,\y+0.18) (3.7,\y-0.08) (5.8,\y+0.14)};
    \node[right] at (5.85,\y+0.14) {\(\lab\)};
  }
  \draw[blue!70!black,very thick,->]
    (0.9,2.02) .. controls (2.0,2.28) and (3.4,1.88) .. (5.1,2.10)
    node[midway,below=4pt]{\(X_H\)};
  \draw[orange!85!black,very thick,->]
    (2.1,0.72) .. controls (2.55,1.55) and (3.0,2.42) .. (3.55,3.18)
    node[pos=0.62,right=4pt]{generic \(X_G\)};
  \draw[green!45!black,very thick,->]
    (1.2,3.23) .. controls (2.3,3.48) and (3.8,3.08) .. (5.0,3.30)
    node[midway,above=4pt]{symmetry flow};
\end{tikzpicture}
}
\caption{Reconstruction of the board geometry: a generic generator crosses
energy surfaces, whereas a symmetry generator is tangent to them.}
\label{fig:l09-flow-reconstruction}
\end{figure}

Equation~\eqref{eq:l09-general-generated-flow} applied to \(A=H\) gives
\begin{equation}
 \delta H=\epsilon\{H,G\}.
\end{equation}
The symmetry condition is therefore
\begin{equation}
 \{H,G\}=0.
 \label{eq:l09-symmetry-condition}
\end{equation}
By antisymmetry, \(\{G,H\}=0\). If \(G\) has no explicit time dependence,
Hamiltonian evolution gives
\begin{equation}
 \dot G=\{G,H\}=0.
 \label{eq:l09-conservation}
\end{equation}

\begin{figure}[tbp]
\centering
\includegraphics[width=0.72\textwidth]{lecture_09_figure_12.png}
\caption{The symmetry condition at 01:05:44:
\(\{H,G\}=0\) is equivalently \(\{G,H\}=0\).}
\label{fig:l09-symmetry-board}
\end{figure}

This is the reciprocal core of Noether's theorem in Hamiltonian language.
The \(G\)-flow does not change \(H\), and the \(H\)-flow does not change
\(G\). A generator of a continuous symmetry is a conserved quantity.

The first statement, that \(G\) leaves the energy unchanged, is
\(\{H,G\}=0\). Antisymmetry gives \(\{G,H\}=0\), and the latter bracket is
precisely \(\dot G\). They are the same vanishing bracket read along the two
flows.

\begin{classroomqa}
\audiencequestion{What is a concrete example of this abstract relation?
\lecturetimestamp{01:07:56--01:08:40}}
\lecturerresponse{Rotational symmetry and angular momentum. For a rotationally
invariant Hamiltonian, the angular-momentum generator has zero Poisson bracket
with \(H\), so angular momentum is conserved.}
\end{classroomqa}

\section{Angular Momentum and the Free Particle}

Take a free particle in two dimensions:
\begin{equation}
 H=\frac{p_x^2+p_y^2}{2m},
 \qquad
 G=L_z=xp_y-yp_x.
 \label{eq:l09-free-angular-pair}
\end{equation}

\begin{figure}[tbp]
\centering
\includegraphics[width=0.74\textwidth]{lecture_09_figure_13.png}
\caption{The free-particle Hamiltonian and angular-momentum generator at
01:09:35. The board sets \(m=1\); the derivation here retains \(m\).}
\label{fig:l09-angular-pair-board}
\end{figure}

Only canonical pairs contribute:
\[
 \{x,p_x\}=1,\qquad \{y,p_y\}=1,
\]
while \(\{x,p_y\}=\{y,p_x\}=0\). Expanding,
\begin{align}
 \{G,H\}
 &=\frac{1}{2m}
 \{xp_y-yp_x,p_x^2+p_y^2\}\notag\\
 &=\frac{1}{2m}\{xp_y,p_x^2\}
  +\frac{1}{2m}\{xp_y,p_y^2\}\notag\\
 &\quad-\frac{1}{2m}\{yp_x,p_x^2\}
  -\frac{1}{2m}\{yp_x,p_y^2\}.
\end{align}
The two cross-mismatched terms vanish. The remaining terms are
\begin{align}
 \{xp_y,p_x^2\}
 &=p_y\{x,p_x^2\}=2p_xp_y,\notag\\
 \{yp_x,p_y^2\}
 &=p_x\{y,p_y^2\}=2p_xp_y.
\end{align}
They cancel:
\begin{equation}
 \{G,H\}
 =\frac{1}{2m}\left(2p_xp_y-2p_xp_y\right)=0.
 \label{eq:l09-angular-cancellation}
\end{equation}

\begin{figure}[tbp]
\centering
\includegraphics[width=0.76\textwidth]{lecture_09_figure_14.png}
\caption{The cancellation at 01:12:50. The two surviving product-rule terms
are equal and enter with opposite signs.}
\label{fig:l09-cancellation-board}
\end{figure}

The same pair can be read in reverse. Invent a system whose Hamiltonian is
the rotation generator,
\begin{equation}
 H_{\mathrm{rot}}=xp_y-yp_x.
\end{equation}

\begin{classroomqa}
\audiencequestion{What would it physically mean to use the generator
\(xp_y-yp_x\) itself as a Hamiltonian?
\lecturetimestamp{01:14:30--01:16:35}}
\lecturerresponse{Work out Hamilton's equations. The generator then produces
uniform rotation: the position and momentum vectors each move around a
circle.}
\end{classroomqa}

Hamilton's equations give
\begin{equation}
 \dot x=-y,\qquad \dot y=x,\qquad
 \dot p_x=-p_y,\qquad \dot p_y=p_x.
 \label{eq:l09-rotation-flow}
\end{equation}
Both the position vector and the momentum vector rotate uniformly. In
particular,
\[
 \ddot x=-x,\qquad \ddot y=-y,
\]
so the configuration-space orbit is a circle.

\begin{figure}[tbp]
\centering
\includegraphics[width=0.74\textwidth]{lecture_09_figure_15.png}
\caption{The circular orbit drawn at 01:16:35 after choosing angular momentum
itself as the Hamiltonian.}
\label{fig:l09-circle-board}
\end{figure}

Because
\(\{p_x^2+p_y^2,H_{\mathrm{rot}}\}=0\), the squared momentum is conserved
in this invented system. The algebra does not privilege either member of a
vanishing-bracket pair; physics tells us which one is the actual Hamiltonian.

\begin{classroomqa}
\audiencequestion{With many degrees of freedom, does each transformed
coordinate depend only on the old pair carrying the same index?
\lecturetimestamp{01:17:01--01:18:09}}
\lecturerresponse{No. There is one transformed \(Q_i\) and \(P_i\) for each
index, but each may be a function of all the old \(q_j\)'s and \(p_j\)'s. The
full set of canonical bracket relations constrains those coupled functions.}
\end{classroomqa}

\begin{classroomqa}
\audiencequestion{How is the momentum related to the velocity in Hamiltonian
language? It is not always \(m\dot q\).
\lecturetimestamp{01:18:12--01:21:08}}
\lecturerresponse{Use the first Hamilton equation,
\(\dot q_i=\partial H/\partial p_i\), and solve it for the momenta when that
relation is invertible. For
\(H=p^2/(2m)+U(q)\), it gives \(p=m\dot q\). The second equation,
\(\dot p=-\partial U/\partial q\), then gives Newton's force law. More general
Hamiltonians produce more general momentum--velocity relations.}
\end{classroomqa}

The bracket that was introduced as an abstract algebra has now done three
jobs. It generated time evolution, characterized the coordinate changes that
preserve mechanics, and converted symmetry into conservation. In quantum
mechanics the corresponding algebra becomes even more concrete: vanishing
Poisson brackets give way to commuting operators, and generators act directly
on states and observables. The classical structure already contains the map.
